% Generated by `taskgen course examples` from number_theory.order_product_tree (seed 3).
% Do not edit: change examples.yaml and regenerate.
\begin{example}[порядок с помощью дерева произведений]\label{ex:order-product-tree}
Определить порядок элемента $37$ поля вычетов $\mathbb{Z}/421\mathbb{Z}$ с помощью построения дерева произведений.

\par\noindent\textit{Решение.} $N = p - 1 = 420 = 2^{2} \cdot 3 \cdot 5 \cdot 7$. Листья дерева --- примарные делители $4, 3, 5, 7$; дерево произведений (от корня к листьям): $420$; $12,\ 35$; $4,\ 3,\ 5,\ 7$.\par Спуск: корню соответствует $h = g = 37$; узел с элементом $h$ передаёт левому потомку $h^{P_R}$, правому $h^{P_L}$ ($P_L, P_R$ --- произведения в поддеревьях), так что лист $q^e$ получает $g^{N/q^e}$:\par\noindent узел $\{4, 3, 5, 7\}$, $h = 37$: $h^{35} \equiv 392$, $h^{12} \equiv 49$;\par\noindent узел $\{4, 3\}$, $h = 392$: $h^{3} \equiv 29$, $h^{4} \equiv 1$;\par\noindent узел $\{5, 7\}$, $h = 49$: $h^{7} \equiv 279$, $h^{5} \equiv 247$;\par Порядок элемента листа $q^e$ равен $q^k$, где $k$ --- первый шаг, на котором последовательность $h, h^q, h^{q^2}, \dots$ достигает $1$:\par\noindent $4$: $29 \to 420 \to 1$, $k = 2$;\par\noindent $3$: $1$, $k = 0$;\par\noindent $5$: $279 \to 1$, $k = 1$;\par\noindent $7$: $247 \to 1$, $k = 1$;\[\ord(37) = 2^{2} \cdot 3^{0} \cdot 5^{1} \cdot 7^{1} = 140.\]
\end{example}
